\chapter{Introduction and Overview}
\label{ch:intro}

\section{Background and Significance}
Biological mechanisms underlying complex traits... As described by Smith et al. \cite{Smith2024}, methylation clocks provide powerful biomarkers of biological age.

\section{Research Objectives}
The overall objective of this dissertation is to:
\begin{enumerate}
    \item Develop generalizable methylation clock models across diverse populations.
    \item Evaluate non-coding variant effects using deep generative genomics.
\end{enumerate}

\section{Figures and Tables Example}
Below is an example of an incorporated figure (\cref{fig:pipeline_overview}).

\begin{figure}[htbp]
    \centering
    \singlespacing
    % \includegraphics[width=0.8\textwidth]{pipeline_overview.pdf}
    \caption[Overview of the dissertation research workflow]{{\bfseries Overview of the dissertation research workflow.} Conceptual illustration of integrating epigenomic datasets, variant predictions, and machine learning models.}
    \label{fig:pipeline_overview}
\end{figure}

\begin{table}[htbp]
    \centering
    \singlespacing
    \caption[Summary of cohorts and sample sizes]{{\bfseries Summary of cohorts and sample sizes across chapters.}}
    \label{tab:sample_summary}
    \begin{tabular}{llrr}
        \toprule
        \textbf{Chapter} & \textbf{Dataset} & \textbf{Samples ($N$)} & \textbf{Tissue Type} \\
        \midrule
        Chapter 2 & Multi-ethnic Epigenetic Cohort & 4,500 & Whole Blood \\
        Chapter 3 & ENCODE / Roadmap Epigenomics & 1,200 & Cell Lines / Primary Tissues \\
        \bottomrule
    \end{tabular}
\end{table}
